\section*{Overview}

The sparse table is the cleanest tool for static range queries on an immutable array. It preprocesses answers on
power-of-two intervals and then recombines those precomputed pieces during queries.

Its best-known use is range minimum query in \(O(1)\) time, but the real lesson is more general: if the array never
changes, it is often worth paying more preprocessing to make every query almost free.

\section*{When to Use It}

Use it when:

\begin{itemize}
  \item the array is static,
  \item you have many queries,
  \item the operation is associative,
  \item and ideally idempotent if you want the \(O(1)\) overlap trick.
\end{itemize}

Range minimum, maximum, and gcd are the classic fits. Range sum does \emph{not} get the same \(O(1)\) trick because
overlapping intervals would double-count values.

\section*{Core Idea}

Precompute answers for every interval of length \(2^k\):
\[
\texttt{st[k][i]} = \text{answer on } [i, i + 2^k - 1].
\]

Then a larger interval can be described using powers of two. For example, length \(13\) can be built from
\(8 + 4 + 1\), which already suggests an \(O(\log n)\) query method.

For idempotent operations such as \(\min\), there is an even better trick: any interval \([l, r]\) can be covered by
two overlapping blocks of equal length \(2^k\), where \(k = \lfloor \log_2(r - l + 1) \rfloor\):
\[
[l, l + 2^k - 1]
\quad\text{and}\quad
[r - 2^k + 1, r].
\]

\section*{Key Insight}

The phrase I keep in mind is: \textbf{static data lets you precompute aggressively}.

The second key fact is the idempotent overlap trick:

\begin{itemize}
  \item for \(\min\), \(\max\), or \(\gcd\), combining two overlapping blocks is safe,
  \item for \(\sum\), the overlap would count some positions twice, so that shortcut breaks.
\end{itemize}

That is why sparse tables feel magical for RMQ but much less magical for sums.

\section*{Operations / Main Technique}

The two main phases are:

\begin{itemize}
  \item \textbf{build}: fill level \(0\) with the array, then build larger powers of two from smaller ones,
  \item \textbf{query}: either decompose the interval into powers of two, or for idempotent operations use the two-block trick.
\end{itemize}

\paragraph{Worked example.}
To answer \(\min\) on a length-\(13\) interval, take \(k = 3\), so the block size is \(8\). Then query the left block
of length \(8\) and the right block of length \(8\), and take the minimum of those two answers. The overlap is harmless
because taking \(\min\) twice does not change the result.

\section*{Correctness Intuition}

Each table entry is correct by construction: it is built from two correct halves of equal length.

For the \(O(1)\) RMQ trick, the queried interval is fully covered by the two chosen blocks. Because the operation is
idempotent, elements that appear in both blocks do not cause trouble:
\[
\min(x, x) = x,
\qquad
\gcd(x, x) = x.
\]

So combining the two block answers gives the same result as combining the entire interval directly.

\section*{Complexity Analysis}

\begin{itemize}
  \item preprocessing: \(O(n \log n)\),
  \item RMQ / idempotent query: \(O(1)\),
  \item associative but non-idempotent query by decomposition: \(O(\log n)\),
  \item memory: \(O(n \log n)\).
\end{itemize}

That memory cost is the main reason not to use sparse tables casually on very large dynamic data.

\section*{Implementation}

The reference code implements the standard RMQ version:

\begin{itemize}
  \item precomputed floor logs,
  \item table levels by powers of two,
  \item \(O(1)\) range minimum query.
\end{itemize}

I like this as a baseline because the min version makes the idempotent requirement completely explicit.

\section*{Common Pitfalls}

\begin{itemize}
  \item Using sparse tables on arrays that change between queries.
  \item Forgetting that the \(O(1)\) query trick assumes idempotence.
  \item Getting the right block start wrong in \(\texttt{st[k][r - (1 << k) + 1]}\).
  \item Building too many columns or using the wrong maximum log.
  \item Choosing a sparse table when a prefix sum or Fenwick tree would be simpler and lighter.
\end{itemize}

\section*{Variants / Extensions}

\begin{itemize}
  \item sparse table for range max or range gcd,
  \item sparse table over any associative operation with \(O(\log n)\) queries,
  \item disjoint sparse table for \(O(1)\) queries on more general associative operations,
  \item RMQ as a subroutine for LCA after Euler tour reduction.
\end{itemize}

The disjoint sparse table is worth knowing exists, even if the ordinary sparse table is still the one I use most.

\section*{Practice Problems}

\begin{itemize}
  \item Static range minimum or gcd queries.
  \item LCA reductions that rely on RMQ over an Euler tour.
  \item Problems with many queries and no updates, where a segment tree would be unnecessary.
\end{itemize}

\section*{References}

\begin{itemize}
  \item \href{https://cp-algorithms.com/data_structures/sparse-table.html}{cp-algorithms: Sparse Table}.
  \item \href{https://usaco.guide/CPH.pdf}{Competitive Programmer's Handbook}.
  \item \href{https://codeforces.com/catalog}{Codeforces Catalog}.
\end{itemize}
